% Build from the repository root; keep auxiliary files outside this folder:
% mkdir -p /tmp/lewm-odessy-pdf
% pdflatex -halt-on-error -interaction=nonstopmode -output-directory=/tmp/lewm-odessy-pdf theOdessy/reproduction.tex
% cp /tmp/lewm-odessy-pdf/reproduction.pdf theOdessy/reproduction.pdf
\documentclass[11pt,a4paper]{article}
\usepackage[margin=23mm]{geometry}
\usepackage[T1]{fontenc}
\usepackage{lmodern}
\usepackage{microtype}
\usepackage{booktabs,tabularx}
\usepackage{xcolor}
\definecolor{accent}{HTML}{23485B}
\usepackage[colorlinks=true,urlcolor=accent,linkcolor=accent]{hyperref}
\hypersetup{
  pdftitle={LeWM checkpoint reproduction},
  pdfauthor={Sachin Mohan},
  pdfsubject={Experiment 1: frozen checkpoint competence on nearby goals}
}
\setlength{\parindent}{0pt}
\setlength{\parskip}{8pt}
\setlength{\emergencystretch}{2em}
\pagestyle{empty}

\begin{document}
{\small\sffamily\color{accent} theOdessy \quad / \quad Experiment 1}

{\LARGE\bfseries LeWM checkpoint reproduction}

{\large Checking the foundation before building on it}
\vspace{3pt}

We want to establish that the released LeWorldModel (LeWM) checkpoints can
complete the tasks reported in the paper using its nearby-goal evaluation.
This gives our later planning and safety research a tested starting point.
The test covers the complete controller: frozen model, image processing,
planner and environment.

We will use the authors' available checkpoint for each environment and keep
its weights fixed. The first comparison uses the original LeWM planner with
the Cross-Entropy Method (CEM). Anchored Planning and our safety adaptations
come after this baseline check; world-model retraining is outside this study.

\begin{center}
\renewcommand{\arraystretch}{1.12}
\begin{tabularx}{\linewidth}{@{}lXr@{}}
\toprule
\textbf{Environment} & \textbf{Task} & \textbf{Paper success} \\
\midrule
TwoRoom & Navigate between rooms through a doorway & 87\% \\
Reacher & Match a target arm configuration & 86\% \\
PushT & Push a T-shaped block into position & 96\% \\
OGBench Cube & Move a single cube to its target & 74\% \\
\bottomrule
\end{tabularx}
\end{center}

For each case, the start and goal come from the same recorded trajectory.
The goal is \textbf{25 environment actions ahead}, and the controller has
\textbf{50 actions} to reach it. CEM uses 300 candidates and 30 elites.
It plans five blocks of five actions and executes the 25-action plan before
replanning. These scores describe nearby-goal control, rather than arbitrary
full-task completion.

We will begin with the released evaluation recipe of 50 cases per environment
and seed 42, saving the exact sampled cases. Additional evaluation seeds will
be chosen beforehand and reported separately. Checkpoint and dataset versions,
preprocessing, action normalization and success rules will be fixed. Before
running, we will resolve the paper/code differences in CEM iteration counts
and TwoRoom history length and record the settings we adopt.

The outcome we want is a repeatable baseline with performance consistent with
the published results. We will report successes, failures and uncertainty,
using an acceptable deviation chosen before seeing the outcomes. Exact equality
with a rounded paper percentage is not required. Any substantial shortfall
will be investigated and recorded before we build on that environment.

The final record will contain a results table, the exact configuration and
case list, episode outcomes, simulator actions, model queries and elapsed time.
This establishes task competence under the tested conditions; safety and
longer-horizon performance require their own evaluations.

\vfill
{\footnotesize
Planning overview, 4 November 2026. Final settings, repeats, acceptance margins
and compute budget will be fixed before execution.
}
\end{document}
